\appendixsection{Допущения, ограничения и проверка полноты рукописи}
\label{app:assumptions}

\subsection{Проверка полноты обязательных частей}

Состав подготовленной рукописи проверен по обязательным элементам ВКР. Результат
сведён в таблицу~\ref{tab:appendix-completeness-check}.

\begin{table}[h]
    \centering
    \footnotesize
    \renewcommand{\arraystretch}{1.15}
    \caption{Проверка полноты обязательных частей ВКР}
    \label{tab:appendix-completeness-check}
    \begin{tabular}{|L{0.34\textwidth}|L{0.14\textwidth}|L{0.38\textwidth}|}
        \hline
        \textbf{Обязательная часть} & \textbf{Статус} & \textbf{Комментарий} \\
        \hline
        Титульный лист & Есть & Формируется средствами шаблона diploma. \\
        \hline
        Реферат & Есть & Подготовлен как отдельный файл главы. \\
        \hline
        Содержание & Есть & Формируется автоматически. \\
        \hline
        Введение & Есть & Содержит актуальность, цель, задачи, объект и предмет. \\
        \hline
        Аналитическая глава & Есть & Включает анализ предметной области и обоснование выбора технологий. \\
        \hline
        Постановка задачи и требования & Есть & Включены в аналитическую главу отдельным подразделом. \\
        \hline
        Проектирование архитектуры & Есть & Описаны сервисы, потоки данных, надёжность и безопасность. \\
        \hline
        Моделирование данных & Есть & Описано в главе реализации и приложении. \\
        \hline
        Описание API/контрактов & Есть & Вынесено в главу реализации и приложение с контрактами сервисного взаимодействия. \\
        \hline
        Реализация & Есть & Отражены реальные сервисы и фактический кодовый пайплайн. \\
        \hline
        Тестирование & Есть & Описаны методика, критерии успешности и артефакты наблюдаемости. \\
        \hline
        Экономическая часть / безопасность & Есть & Выполнена качественная оценка и анализ рисков. \\
        \hline
        Заключение & Есть & Сформулированы результаты и направления развития. \\
        \hline
        Список источников & Есть & Оформляется через biblatex. \\
        \hline
        Приложения & Есть & Добавлены приложения по контрактам, Kafka, данным и наблюдаемости. \\
        \hline
        Задание на ВКР & Требует внешнего документа & См. допущение ниже. \\
        \hline
        Календарный план & Требует внешнего документа & См. допущение ниже. \\
        \hline
    \end{tabular}
\end{table}

\subsection{Критические допущения и ограничения}

\assumption{В текущий LaTeX-проект не включены отдельные файлы задания на ВКР и
календарного плана. Согласно используемому шаблону института эти документы
оформляются как внешние материалы и должны быть приложены к итоговому комплекту
работы отдельно.}

\assumption{Текущая кодовая база реализует полноценные sender-сервисы только для
каналов электронной почты и SMS. Канал push присутствует в модели профиля
получателя и может быть развит в дальнейшем, но отдельный push-sender в проекте
на момент подготовки отчёта отсутствует.}

\assumption{В демонстрационном профиле текущей реализации часть gRPC-сервисов и
клиентов работает без обязательной аутентификации, а Kafka и отдельные gRPC-клиенты
используют plaintext-транспорт. Для промышленного внедрения требуется включить
TLS, сервисную аутентификацию и централизованное управление секретами.}

\assumption{Часть адаптеров доставки может работать в тестовом режиме log-only
или через логирующий gateway. Это позволяет верифицировать программный пайплайн,
но не заменяет отдельную промышленную проверку реальных внешних провайдеров.}

\assumption{Экономическая глава выполнена как качественная оценка, поскольку в
составе репозитория отсутствуют утверждённые числовые исходные данные, необходимые
для полного нормативного калькуляционного расчёта.}

\subsection{Проверка целостности текста}

При доработке рукописи были дополнительно проверены следующие аспекты:
\begin{enumerate}
    \item единообразие терминов событие, запуск доставки, канальная доставка;
    \item отсутствие в основном тексте утверждений о несуществующем push-sender;
    \item согласованность описания data flow с фактическими gRPC- и Kafka-контрактами;
    \item наличие источников для обзорных и теоретических положений;
    \item наличие выводов по каждой основной главе.
\end{enumerate}

